Em primeiro lugar, gostaria de agradecer ao Professor Nuno Gonçalves, por sempre me ter acompanhado durante a licenciatura, mestrado, investigação, por me ter aceite o meu pedido de orientador de dissertação e, por ter persistido comigo durante estes últimos dois anos, sem nunca desistir de mim. Muito obrigado por toda a sua orientação, atenção, paciencia, horas de reuniões e pela sua amizade. Levá-lo-ei no coração para a vida. Ao meu co-orientador Professor João Marcos, que apesar de teimar que o trate por "tu", acompanhou-me já em mestrado, causando uma boa dinâmica de trabalho e amizade. Pelos ensinamentos, dentro e fora, da sala de aula, e pela partilha de experiências, sou-lhe muito grato.

A ti Coimbra, Universidade de Coimbra e DEEC, por me monirem das condições de estudo e de vida, pelas amizades, vivências, e aprendizagens que, orgulhosamente, me moldaram na pessoa que hoje sou. Obrigado a todos os docentes que contribuíram para a minha formação e crescimento pessoal, em particular ao Professor Diogo Júdice, Gabriel Falcão, Paulo Coimbra, e Lúcia Martins.

Professor César, Professora Adélia, consegui. Os tempos de escola já lá vão, mas não me esqueci de vocês. Tinha de reservar um cantinho para os meus professores do ensino secundário, A si, Professora Adélia, por me ter dado na cabeça quando precisava, e a si Professor César, por me ter sido um porto seguro mesmo após tantos anos sem nos vermos. Obrigado por tudo.

A todos os meus colegas do ISR, pelos ensinamentos e horas de debugging. Obrigado Farhad, Luiz, Aniana e XXX por terem codificados as imagens usads nesta dissertação. E ao VisTeam, pelo o apoio triplicado na avaliação das imagens, se a métrica está com a qualidade que está, certamente é graças a vocês. Um obrigado especial ao Allan, pela paciencia, e ao Diogo, pelo passinho de dança no NL.

NEEEC, meu querido NEEEC, já disse, quando perdi uma casa, ganhei outra. Quando entrei para o curso, não estava à espera de encontrar uma segunda família, mas encontrei-a em vocês. Desde os convívios, aos jogos infindáveis de ping-pong, desde as milhentas horas de "estudo", às canecas no Reis. Obrigado por me terem ensinado que o DEEC é muito mais que um departamento, é um lar, onde devemos reivindicar pela nossa casa, e ajudar o colega do lado.

Flórido, Casca, Zé, eina... acho que não preciso de dizer muito mais... Se o mandato foi curto, a amizade que construímos é para a vida. Eu sei que houve reuniões difíceis, "Onde é que estão os chocolates?", "Flórido, eu meti a mão na urna!", "Casca, e essa gala, vai dar lucro?", "Desculpa Zé, esqueci-me de fazer a ata". Mas nem todas as reuniões foram más, as do Reis corriam bem (pelo que me lembro). Obrigado por me terem convidado neste projeto. Sou-vos muito grato.

Déh e Duarte, nós juntos somos a tripla maravilha. Os três vices que sabem o que é a verdadeira luta. Obrigado pelo carinho especial, horas de conversa, boleias, e apoio emocional, se me disserem que gostam tanto de mim como eu gosto de vocês, estariam a mentir! Amo-vos do fundo do coração.

Familia de praxe, (ali até à Marysol, exclusive). Paulino, seu cabeludo, lá sabia eu que, mesmo a 2000km de distância, me ias fazer chorar! Obrigado pelas turras na cabeça (excepto no Bot), gosto muito de ti e espero que a tua vida te sorria tanto como tu me fazes sorrir. Sei que vais ser um grande Engenheiro, tal como o Padrinho, obrigado por me teres escolhido (vamos ignorar o Casca). Marysol, estava a brincar! Tu desafias a paciencia de qualquer um, espero ter tido um impacto positivo na tua vida, as minhas pancadinhas sempre foram com carinho.

A todos os que fizeram parte do (melhor) Carro da Queima, um especial obrigado ao meu amigo Carlos Faria, João Domigues, Gonçalo Santos, João Magano, e Gonçalo Dias, pela dedicação e acompanhamento, todos praticamente convertidos para informática ou para o Algarve, mesmo longe, trago-vos no coração. Obrigado por me terem proporcionado um dos melhores dias da minha vida.

Aos Amontoados, a ti Andreia, Tiago, Eduardo, João, Bernardino e Catarina, por terem feito parte da minha integração e, da minha vida. Eduardo, ainda estou à espera que faças a laboratorial de SM! Obrigado por teres estado lá nos melhores, e piores momentos. Tiago, és uma força da natureza, o Iron Man está ao virar da esquina, não tenho e menor dúvida que te vais tornar um. Para mim serás sempre o meu irmão, quando cortarem a fita, berra e grita, na tua vida só tu ditas o guião.

Miguel Álvila, Daniel Agostinho, tinha de haver um paragrafo vosso. A palavra amizade não chega para descrever o que existe entre nós. Desde aos assaltos às bancas de shots no 24, às noites de boémia ao som do GrandSons; do Reis até Vigo, de Anadia até Viana. A tripla não dorme, nem deixa que Coimbra inteira adormeça. Obrigado por terem estado do meu lado este último ano, sei que isto é só um começo. Um brinde a vocês, sei que brevemente estaremos juntos a celebrar da mesma maneira de sempre. Jaci, Swarovski, não me esqueci de vocês, mas o parágrafo já estava grande demais. Um beijinho grande às duas, obrigado por me aturarem e desculpem lá qualquer coisa!

Augusto, Cunha, Vasco, Pedro, para quando um pedi-tascas? Pedro, ainda te continuas a atrasar prá escolinha? Cunha, os 20 aqui são mais dificeis. Vasco, quando é que voltas a Portugal? Augusto, já entras-te em direito? Tinha de desabafar... Obrigado por terem feito parte da minha vida, estarei sempre aqui para vocês. Danny, my man, if I'd wrote in Portuguese you wouldn't understand a word besides "bumbum tantan". I know you are far away, It's about time I'd come to visit you now. It's only fair. Undoubtedly one of the most unlikely friendship one could ever make, you are the sole reason I believe in destiny, love you bro.

A ti, Inês, por me teres aturado, por me teres ajudado a crescer, por me teres ensinado tanto, e por me teres feito tão feliz. Obrigado por tudo o que és e por tudo o que me fazes ser. Amo-te.

Mãe, pai, nuna. Mãe, obrigado por tudo, sei que ainda não me tornei no Grande Homem que tu tanto querias, mas estou a trabalhar para isso, quem tem uma mãe, tem tudo, mas quem tem uma mãe como tu, não pode pedir por mais nada na vida, amo-te, incondicionalmente. Pai, sei que nem sempre foi fácil, sei que nestes últimos 6 anos, passei muito pouco tempo contigo, obrigado pelo carinho, e por me teres deixado o bichinho pela política, do teu filhote réré, amo-te. Nuninha, minha maninha, este é um momento especial e não te podia deixar de fora, temos muito mais em comum do que tu pensas, o mano ama-te, mesmo que sejas parvinha às vezes.

Deonilde, tinha de guardar o melhor para o fim. Espero um dia vir a ser tudo aquilo que dizes que sou... mas por agora, espero que estejas orgulhoso do teu Neto. Ainda me lembro, de te ligar, e te contar quantas frequências e trabalhos tinha para a próxima semana, e, tu, na tua inocência, respondeste "O teu trabalho é estudar... não te pedem mais nada, pois não?" Nunca deixes de rezar por mim.
